%%%%%%%%%%%%%%%%%%%%%%%%%%%%%%%%%%%%%%%%%
% Medium Length Professional CV
% LaTeX Template
% Version 2.0 (8/5/13)
%
% This template has been downloaded from:
% http://www.LaTeXTemplates.com
%
% Original author:
% Trey Hunner (http://www.treyhunner.com/)
%
% Important note:
% This template requires the resume.cls file to be in the same directory as the
% .tex file. The resume.cls file provides the resume style used for structuring the
% document.
%
%%%%%%%%%%%%%%%%%%%%%%%%%%%%%%%%%%%%%%%%%%%%%%%%%%%%%%%%%%%%%%%%%%%%%%%%%%%%%%%%%%
% Medium Length Professional CV
% LaTeX Template
% Version 2.0 (8/5/13)
%
% This template has been downloaded from:
% http://www.LaTeXTemplates.com
%
% Original author:
% Trey Hunner (http://www.treyhunner.com/)
%
% Important note:
% This template requires the resume.cls file to be in the same directory as the
% .tex file. The resume.cls file provides the resume style used for structuring the
% document.
%
%%%%%%%%%%%%%%%%%%%%%%%%%%%%%%%%%%%%%%%%%

%----------------------------------------------------------------------------------------
%	PACKAGES AND OTHER DOCUMENT CONFIGURATIONS
%----------------------------------------------------------------------------------------

\documentclass{resume} % Use the custom resume.cls style

\usepackage{hyperref}
\usepackage[left=0.45in,top=0.35in,right=0.45in,bottom=0.35in]{geometry} % Document margins

  % If using pdflatex:
\usepackage[utf8]{inputenc}
\usepackage[T1]{fontenc}
\usepackage[default]{lato}
\usepackage{fontawesome}
\usepackage{setspace}
\usepackage{biblatex}
\usepackage{dashrule}
\usepackage{graphicx}
\usepackage{tikz}
\usepackage{hyperref}
\hypersetup{
    colorlinks=true,
    linkcolor=blue,
    filecolor=magenta,
    urlcolor=cyan,
}

\newcommand*\googlescholaricon{\raisebox{-0.3em}{\includegraphics[width=1em]{google-scholar.pdf}}}

\newcommand{\tab}[1]{\hspace{.2667\textwidth}\rlap{#1}}
\newcommand{\itab}[1]{\hspace{0em}\rlap{#1}}
\name{Shreyas Kalvankar} % Your name
\address{\faMapMarker ~ Nashik, Maharashta, India} % Your address
%\address{123 Pleasant Lane \\ City, State 12345} % Your secondary addess (optional)
\address{\footnotesize \faPhone ~ (+91)~9423555723 ~ \faEnvelope ~ shreyaskalvankar@gmail.com ~ \faLinkedin ~ \href{https://linkedin.com/in/shreyas-kalvankar}{linkedin.com/in/shreyas-kalvankar} ~ \faGithub ~ \href{https://github.com/obi-wan-shinobi}{github.com/obi-wan-shinobi} ~ 
\googlescholaricon ~ \href{https://scholar.google.com/citations?hl=en&user=W8-T__UAAAAJ}{Scholar}}

\addbibresource{references.bib}

\begin{document}

%----------------------------------------------------------------------------------------
%	EDUCATION SECTION
%----------------------------------------------------------------------------------------
\smallskip
\begin{rSection}{Education}

{\bf Bachelor of Engineering (Computer Engineering)} \hfill {\em 2017 - 2021}
\\ K.K. Wagh Institute of Engineering  \hfill { Overall GPA: 9.7/10}
\\ Education \& Research, Nashik    \hfill {\small{(Rank 1)}}

\smallskip

% {\bf Higher Secondary Certificate} \hfill {\em 2017}
% \\ HPT Arts \& RYK Science College, Nashik \hfill { Percentage: 87.07\%}

\end{rSection}
\begin{rSection}{PUBLICATIONS}
\vspace{-1em}
\item Shreyas Bapat et al. \textit{EinsteinPy: A Community Python Package for General Relativity.} 2020.\\ arXiv: \href{https://arxiv.org/abs/2005.11288}{2005.11288 [gr-qc]}.

\item Shreyas Kalvankar et al. \textit{Galaxy Morphology Classification using EfficientNet Architectures.} 2020.\\ arXiv:
\href{https://arxiv.org/abs/2008.13611}{2005.13611 [cs.CV]}.

\item Kalvankar, Sh., Pandit, Hr., Parwate, Pr., Patil, At. \& Kamalapur, Sn., (2022). \textit{Astronomical Image Colorization and Up-scaling with Conditional Generative Adversarial Networks.}\\ In: Demmler, D., Krupka, D. \& Federrath, H. (Hrsg.), INFORMATIK 2022. Gesellschaft für Informatik, Bonn. (S. 489\-498).
\href{https://dl.gi.de/bitstream/handle/20.500.12116/39540/mlastro\_04.pdf}{DOI: 10.18420/inf2022\_40}.
\end{rSection}

\smallskip
%----------------------------------------------------------------------------------------
%	WORK EXPERIENCE SECTION
%----------------------------------------------------------------------------------------
\smallskip
\begin{rSection}{Professional Experience}

\begin{rSubsection}{Dalton Maag Ltd.}{November 2021 - Present}{}{}
\begin{rSubsubsection}{\textbf{Software Developer}}{London, United Kingdom}{}
\item Developed a system for automatically generate Chinese-Japanese-Korean font glyphs using genetic algorithms.
\item Implemented glyph data model for Devanagari script from scratch and improved data models for Arabic, Greek and Cyrillic scripts in the internal pricing application, enabling accurate pricing for non-Latin projects.
\item Upgraded company's pricing model in the internal pricing application to provide faster, accurate quotes, boosting efficiency in project planning and increasing productivity.
\end{rSubsubsection}
\vspace{-0.4em}
\end{rSubsection}
\vspace{-0.4em}
\hdashrule[1pt]{19cm}{0.5pt}{1mm 1mm}
\begin{rSubsection}{Relfor Labs Pvt. Ltd.}
{Pune, India}{}{}
\begin{rSubsubsection}{\textbf{Machine Learning Research Scientist}}{September 2022 - Present}{}
    \item Created novel neural architectures which out-performed the previous best-performing models, on additional data, by \textbf{2\%}, achieving an accuracy of \textbf{98.6\%} in audio data classification.
    \item Set up the ML training pipeline using PyTorch lightning on Nvidia DGX A100, with automatic hyperparameter tuning using Optuna and dynamic architecture updates, expediting experimentation.
    \item Experimented with several data models and improved the training pipeline, making experimentation \textbf{$\sim$10\%} faster, which boosted performance metrics by \textbf{$\sim$1\%}
\end{rSubsubsection}\vspace{-1em}
\begin{rSubsubsection}{\textbf{Machine Learning Engineer}}{August 2021 - November 2021}{}
    \item Designed novel deep Convolutional Neural Network architectures for audio data classification which beat state-of-the-art models.
    \item Performed extensive research on the designed architectures \& network embeddings and created \textbf{novel loss functions} and a \textbf{custom gated unit block} achieving \textbf{>98\%} accuracy and \textbf{$\sim$0.98} F1-score.
    \item Implemented various statistical methods to boost model performance by analyzing threshold values and increasing precision to \textbf{$\sim$98\%} while maintaining high accuracy \textbf{>98\%}.
\end{rSubsubsection}
\vspace{-0.4em}
\end{rSubsection}
\vspace{-0.4em}
\hdashrule[1pt]{19cm}{0.5pt}{1mm 1mm}
\begin{rSubsection}{FinIQ Consulting India Pvt. Ltd.}{May 2020 - June 2020}{}{}
\begin{rSubsubsection}{\textbf{Software Development Intern}}{Nashik, India}{}
\item Developed a front-end using AngularJS for forex trading with interactive visualization and chatbot service, providing an appealing platform for forex operations.
\item Created a python module for stress testing CPU and memory as per user input using variable load calibration.
\item GitHub: \href{https://github.com/obi-wan-shinobi/Stress-test}{CPU and Memory Stressing module} \& \href{https://github.com/obi-wan-shinobi/Forex-Trading}{Forex Trading Platform}.
\end{rSubsubsection}
\end{rSubsection}


\end{rSection}
\smallskip
%	EXAMPLE SECTION
%----------------------------------------------------------------------------------------

%----------------------------------------------------------------------------------------

%----------------------------------------------------------------------------------------
\begin{rSection}{PROJECTS \& RESEARCH}

\begin{rSubsection}{\href{https://github.com/obi-wan-shinobi/Astronomical-Image-Colorization}{Astronomical Image Colorization and Super-resolution using GANs}}{August 2020 - June 2021}{}{}
\vspace{-0.9em}
\item Led a team of four members in a project for automatic colorizing and upscaling low-resolution, grayscale astronomical images.
\item Created a dataset of $\sim 5000$ images by scraping the Hubble archives
\item Developed variations of GAN architectures, effectively creating a novel training method, for colorizing images achieving visually pleasing results
\item Implemented a variation of SRGAN architecture suitable for the data and obtained high-resolution images
\end{rSubsection}

\begin{rSubsection}{\href{https://github.com/obi-wan-shinobi/GalaxyEfficientNets}{The Galaxy Zoo Project}}{August 2019 - September 2020}{}{}
\vspace{-0.9em}
\item A galaxy morphology classification project, based on Kaggle Galaxy Zoo 2 competition, implementing the EfficientNet architectures
\item Developed a CNN for vote fraction predictions of 37 galaxy features from the Galaxy Zoo decision tree with a RMSE score of \textbf{0.07765}, ranking us in the \textbf{top 3} on the public leaderboard.
\item Developed a CNN for classification of galaxies into 7 classes based on their morphologies with an accuracy of \textbf{93.7\%} and an F1 score of \textbf{0.8857}.
\end{rSubsection}

\begin{rSubsection}{\href{https://github.com/einsteinpy/einsteinpy}{The EinsteinPy Project}}{March 2020 - April 2020}{}{}
\vspace{-0.9em}
\item Contributor to an open source community python package for general relativity (\textbf{500+ stars} on GitHub).
\item Added Reissner–Nordström metric: a static solution to the Einstein-Maxwell field equations, into the code.
\item Corrections in the Kerr-Newman and Kerr metrics classes.
\item Added calculations of the event horizon and ergosphere for a Kerr-Newman black hole.
\item \href{https://doi.org/10.5281/zenodo.4445219
}{DOI: 10.5281/zenodo.4445219}
\end{rSubsection}

\begin{rSubsection}{Robocon}{August 2018 - May 2020}{}{}
\vspace{-0.9em}
\item Built a quadruped robot with gait similar to a horse, and a wheeled robot with dynamic locomotive abilities for ABU Robocon 2019
\item Mentored junior members of the team for Robocon 2020; planned and assisted in creating two wheeled robots capable of performing intricate tasks of catching and throwing a rugby ball
\end{rSubsection}

\end{rSection}

\begin{rSection}{Scholastic \& Co-curricular Achievements}
\vspace{-1em}
\item Received the Best Outgoing Student Award in 2021 by the Head of Computer Engineering Dept., K. K. Wagh Institute
\item Received the Award of Academic Excellence for the best academic performance across the institution among 1200 students in all Engg. departments
\item Ranked 9 in phase 1 of ABU Robocon 2019 among 200+ teams across the country
\end{rSection}

%----------------------------------------------------------------------------------------
%	TECHNICAL STRENGTHS SECTION
%----------------------------------------------------------------------------------------
\smallskip
\begin{rSection}{Technical Strengths}

\begin{tabular}{ @{} >{\bfseries}l @{\hspace{6ex}} l }
Computer Languages &  C/C++, Python, Ruby, Javascript, Typescript \\
Web Development & AngularJS, VueJS, ElectronJS, Flask, Ruby on Rails, HTML, CSS\\
Deep Learning Frameworks & Keras, TensorFlow, PyTorch \\
Machine Learning Frameworks & Octave, Sci-kit\\
Embedded Systems & Arduino, RaspberryPi, Teensy \\
Version Control & Git, GitHub \\
Tools & Numpy, Pandas, Scipy, \LaTeX

\end{tabular}

\end{rSection}
%----------------------------------------------------

\end{document}    